%%%%%     PACKS     %%%%%
\documentclass[12pt]{article}
\usepackage[margin=1in,headsep=.60in]{geometry}
\usepackage[utf8]{inputenc}
\usepackage[table, dvipsnames]{xcolor}
\usepackage{array}
\usepackage{amsmath}
\usepackage{booktabs}
\usepackage{mdframed} %For box around text
\usepackage{hyperref} %For hyperlinks in the PDF
\hypersetup{
    colorlinks=true,
    linkcolor=blue,
    filecolor=magenta,
    urlcolor=cyan,
}
\urlstyle{same}
\usepackage{amsthm}
\usepackage{amssymb}
\usepackage{amsfonts}
\usepackage{siunitx}
\usepackage{graphicx}
\usepackage{caption}
\usepackage{pgfplots}
\graphicspath{{Images/}}
\usepackage[colorinlistoftodos]{todonotes}
\usepackage{cleveref}
\usepackage[labelformat=simple]{subcaption}
\usepackage{grffile}
\usepackage{gensymb}
\usepackage{float}
\usepackage[shortlabels]{enumitem}
\usepackage{enumitem}
\setlistdepth{9}
\usepackage{biblatex} %Imports biblatex package
\addbibresource{references.bib}
\usepackage{outlines}
\usepackage{minted}
\usepackage{pdflscape}
\usepackage{everypage}
\usepackage{multirow}
\usepackage{multicol}
\usepackage{afterpage}
\usepackage[labelfont=bf, font=small]{caption}
\usepackage{upgreek}
\usepackage{tikz}
\usepackage{xcolor}
\usepackage{listings}

% style sheet for code
\definecolor{codegreen}{rgb}{0,0.6,0}
\definecolor{codegray}{rgb}{0.5,0.5,0.5}
\definecolor{codepurple}{rgb}{0.58,0,0.82}
\definecolor{backcolour}{rgb}{0.95,0.95,0.92}

\lstdefinestyle{mystyle}{
    backgroundcolor=\color{backcolour},   
    commentstyle=\color{codegreen},
    keywordstyle=\color{magenta},
    numberstyle=\tiny\color{codegray},
    stringstyle=\color{codepurple},
    basicstyle=\ttfamily\footnotesize,
    breakatwhitespace=false,         
    breaklines=true,                 
    captionpos=b,                    
    keepspaces=true,                 
    numbers=left,                    
    numbersep=5pt,                  
    showspaces=false,                
    showstringspaces=false,
    showtabs=false,                  
    tabsize=2
}

\lstset{style=mystyle}

\usepackage{lastpage}
%%%%%     COMMANDS     %%%%%
%%%% Blue box for subsection text (not figures) %%%%
\newenvironment{bluebox}
  {\begin{mdframed}[backgroundcolor=blue!5,linecolor=blue!40,roundcorner=8pt]}
  {\end{mdframed}}

%%%% Simple neutral box for figures %%%%
\newenvironment{figbox}
  {\begin{mdframed}[roundcorner=8pt,shadow=true,shadowsize=4pt,shadowcolor=black!40]}
  {\end{mdframed}}

\begin{document}

\begin{titlepage}
\newcommand{\HRule}{\rule{\linewidth}{0.5mm}}

\center

\textsc{\LARGE Aarhus University}\\[1.5cm]
\textsc{\Large Department Of Mathematics}\\[0.5cm]
\textsc{\large Introduction to Algebra 1}\\[0.5cm]
    

\HRule\\[0.4cm]
    \center 
    {\huge\bfseries  Handin 6}\\[0.4cm] % Title of your document
\HRule\\[1.5cm]

\begin{minipage}{0.4\textwidth}
        \begin{flushleft}
            \large
            \textit{Author}\\   
            Søren M. \textsc{Damsgaard}\\
                 % Your name
        \end{flushleft}
    \end{minipage}
~
    \begin{minipage}{0.4\textwidth}
        \begin{flushright}
            \large
            \textit{Student number}\\
            \textbf{202309814}\\
                 % Studienummer\
            
        \end{flushright}
    \end{minipage}

\vfill\vfill\vfill % Position the date 3/4 down the remaining page
    
    {\large\today}

\vfill\vfill
    \includegraphics[width=0.2\textwidth]{Aarhus_University_seal.png}\\[1cm] % Include a department/university logo - this will require the graphicx package
    

\vfill
\end{titlepage}
%%%%%      CHAPTERS      %%%%%
\hspace{0.02cm}
\section*{Goofy ah Math Number 6}

\section*{RSA-Encryption}
Let Alice have the primes $p = 107$, $q = 127$ and the exponent $k = 5$.
\subsection*{Determine $m$ and $\Phi(m)$ and show that $k = 5$ is a valid choice by computing $\gcd(k, \Phi(m))$.}
from chapter 16.1~\cite{Algebra} it follows that $m = p\cdot q = 107 \cdot 127 = 13589$.
Then, $\Phi(m) = (p-1)(q-1) = 106 \cdot 126 = 13356$.\\
Computing  $\gcd(k, \Phi(m)) = \gcd(5, 13356)$, since $k$ and $\Phi(m)$ needs to be coprime for $k$ to be a valid choice.\\
This can be computed using the Euclidean algorithm 3.2~\cite{Algebra}:
let $r_0 = 13356$, $r_1 = 5$, and compute the sequence of remainders:
\begin{align*}
13356 &= q_1 \cdot 5 + 1 \quad \text{for } q_1 = 2671 \text{ and } r_2 = 1\\ 
5 &= q_2 \cdot 1 + 0 \quad \text{for } q_2 = 5 \text{ and } r_3 = 0
\end{align*}
hence 
\begin{equation*}
    \gcd(5, 13356) = r_2 = 1
\end{equation*}
since the greatest common divisor is 1, according to definition 3.1~\cite{Algebra} then $k = 5$ is coprime to $\Phi(m) = 13356$ and is a certified valid choice for a public key.
\subsection*{Determine Alice's private exponent $u$}
A number $u$ is a valid private exponent if it satisfies the equation $k \cdot u \equiv 1 \mod \Phi(m)$.\\
\newline
To determine this, see definition 9.18~\cite{Algebra} to see how to determine such a number $u$.\\
Some prerequisites to use the method in definition 9.18~\cite{Algebra} is to know the value of $\Phi(m)$ for $m = p\cdot q, \quad p \neq q $ where $p$ and $q$ are primes, and that $k$ is coprime to $\Phi(m)$. All of these prerequisites are satisfied in the previous section.\\
\newline 
The method uses the extended Euclidean algorithm 3.3~\cite{Algebra} to find bezout coefficients $u$ and $v$ such that $k \cdot u + \Phi(m) \cdot v = 1$.\\
Use the algorithm to find $u$:\\
let $ku + \Phi(m)v = \gcd(k, \Phi(m))$, where $k = 5$ and $\Phi(m) = 13356$. the quotients $q_1 = 2671$ and $q_2 = 5$ are already computed in the previous section.\\
\begin{table}[H]
\centering
\begin{tabular}{|c|c|c|c|c|}
\hline
$i$ & $r_i$ & $q_i$ & $u_i$ & $v_i$ \\ \hline
0 & 13356 & -- & 0 & 1 \\ \hline
1 & 5 & 2671 & 1 & 0 \\ \hline
2 & 1 & 5 & -2671 & 1 \\ \hline
\end{tabular}
\vspace{0.2cm}
\caption{Extended Euclidean algorithm for $5$ and $\Phi(m)=13356$}
\label{tab:extended_euclid_rsa}
\end{table}
sanity check:
\begin{align*}
    k \cdot u + \Phi(m) \cdot v = 1\\
    5 \cdot (-2671) + 13356 \cdot 1 = 1\\
\end{align*}
even though the current $u = -2671$ is a valid solution, in accordance to definition 9.18~\cite{Algebra} it needs $u > 0$ and $v < 0$.\\
To accomodate this, let $v' = v \cdot -1 = -1$ and $u' = \Phi(m) - u = 10685 = 13356 -2671$, and then double check that it's still valid:
\begin{align*}
    k \cdot u' + \Phi(m) \cdot v' = 1\\
    5 \cdot 10685 + 13356 \cdot (-1) = 1\\
\end{align*}
which it still is, hence $u' = 10685$ and $v' = -1$ is a valid solution.\\
Then it follow from equation 9.13~\cite{Algebra}:
\begin{equation*}
    ku' + \Phi(m)v' = 1 \implies ku' = 1 - \Phi(m)v' \implies ku' = 1 + \Phi(m)w
\end{equation*}
for $w=-v'$ and letting the multiplier be a modulo $\Phi(m)$, then it follows that
\begin{equation*}
    ku' = 1 + \mod \Phi(m)w \implies ku' \equiv 1 \mod \Phi(m)
\end{equation*}
and for $u' = 10685$ and the rest of the values
\begin{equation*}
    5 \cdot 10685 \equiv 1 \mod 13356
\end{equation*}
and since $u'$ was found via the extended Euclidean algorithm, then $u' = 10685$ is a valid solution to the equation $ku' \equiv 1 \mod \Phi(m)$, which means that $u' = 10685$ is a valid private exponent.\\
\printbibliography
\end{document}